\chapter{Algoritmo Genético para Otimização de \texorpdfstring{$L_0$}{L0}}
\label{ap:ag}

Este apêndice formaliza o algoritmo genético (AG) usado para estimar o
envelope de desempenho da composição do conjunto inicial $L_0$: os limites
superior e inferior que uma composição muito boa, ou muito ruim, consegue
alcançar. Seguem a formulação, os operadores com seus parâmetros e sua
proveniência, e o laço completo, resumido na Figura~\ref{fig:ap-ag-laco}.

\section{Formulação}

Cada \textbf{indivíduo} é um conjunto de $I$ índices únicos do
\textit{pool}, ou seja, um $L_0$ candidato de tamanho fixo. A
\textbf{aptidão} de um indivíduo é o desempenho (Acurácia ou Macro F1) do
classificador leve PVBin treinado nesse candidato. A aptidão é medida em
uma partição de \textbf{aferição}, reservada para este fim e disjunta do
conjunto de teste, porque
otimizar e reportar na mesma partição superajustaria o envelope ao conjunto
de avaliação. Por isso, ao final de cada cenário, o melhor
indivíduo é \textbf{reavaliado no conjunto de teste intocado}, e esse valor,
tipicamente inferior ao de aptidão, é o que se reporta como envelope.

São otimizados \textbf{quatro cenários} por tamanho (maximização e
minimização de cada uma das duas métricas), para dez tamanhos de $L_0$
entre 10 e 30.000: o par de extremos delimita o envelope empírico de
desempenho alcançável pela composição.

\section{Operadores e parâmetros}

Salvo indicação em contrário, os parâmetros abaixo foram definidos no
\textit{notebook} que executou as corridas, que sempre os passa ao
algoritmo (a configuração versionada fixa apenas o tamanho de $L_0$).
População e gerações foram, além disso, conferidas contra os artefatos das
corridas; a população é o único parâmetro sem fonte de configuração, e o
valor reportado é o lido do artefato.

A Tabela~\ref{tab:ap-ag-params} resume os operadores e seus parâmetros.

\begin{table}[htb]
\centering
\caption{Operadores e parâmetros do algoritmo genético.}
\label{tab:ap-ag-params}
\small
\begin{tabular}{p{0.2\textwidth} p{0.24\textwidth} p{0.44\textwidth}}
\toprule
\textbf{Componente} & \textbf{Valor} & \textbf{Observação} \\
\midrule
População e gerações & $N_{pop}=20$; 100 gerações & 2.000 avaliações
supervisionadas por cenário \\
Seleção & torneio, $k_t=3$ & \\
Cruzamento & um ponto, $p_c=0{,}8$ & padrão da implementação é $0{,}7$,
sempre sobrescrito na execução; segue \emph{reparo de unicidade}: índices
repetidos substituídos por sorteio até restaurar $|L_0|=I$ \\
Mutação & $p_m=0{,}1$ & substitui $m_s=\max(1,\lceil 0{,}01 \cdot I\rceil)$
genes por índices inéditos \\
Elitismo & $N_{elite}=2$ & $10\%$ da população \\
\bottomrule
\end{tabular}
\end{table}

O caso $|L_0|=10$ é o único que difere da configuração da tabela: corre por
200 gerações (4.000 avaliações). Nos resultados, esse caso é reportado, como
os demais, na 100\textsuperscript{a} geração ($18{,}82\%$); a corrida segue
até a 200\textsuperscript{a}, em que o melhor indivíduo alcança
$19{,}20\%$. A reexecução independente do envelope, feita para checar a
sensibilidade do resultado à própria busca, usa configuração reduzida:
$N_{pop}=30$ por 40 gerações.

\section{Pseudocódigo}

\begin{algorithm}[H]
  \caption{Estimação do envelope de $L_0$ por algoritmo genético}
  \label{alg:ap-ag}
  \begin{algorithmic}[1]
    \State \textbf{Inicialização:} gerar $N_{pop}$ indivíduos por amostragem
    uniforme sem reposição
    \State Avaliar a aptidão de todos na partição de aferição
    \State \textbf{Laço principal:} por geração, repetir:
    \State \quad copiar a elite ($N_{elite}$ indivíduos)
    \State \quad completar a população: torneio $\to$ cruzamento $\to$
    reparo $\to$ mutação
    \State \quad reavaliar a aptidão na partição de aferição
    \State \textbf{Finalização:} reavaliar o melhor indivíduo no conjunto de
    teste intocado; o valor reportado como envelope é este, não a aptidão
  \end{algorithmic}
\end{algorithm}

\begin{figure}[htbp]
    \centering
    % ============================================================================
% ESQUEMA PROPOSTO (banca, sugestão — não mergeado): o envelope por algoritmo
% genético (Seção sec:metodo-l0-ag) como laço evolutivo + saída com
% anti-circularidade. A prosa da seção carrega o laço inteiro em um
% parágrafo; a figura o desdobra: população -> torneio -> cruzamento ->
% mutação -> aptidão (partição de aferição) -> elitismo -> repete; ao fim,
% reavaliação no teste intocado = o envelope reportado.
%
% Uso no Cap. 3 (dentro de um ambiente figure):
%   % ============================================================================
% ESQUEMA PROPOSTO (banca, sugestão — não mergeado): o envelope por algoritmo
% genético (Seção sec:metodo-l0-ag) como laço evolutivo + saída com
% anti-circularidade. A prosa da seção carrega o laço inteiro em um
% parágrafo; a figura o desdobra: população -> torneio -> cruzamento ->
% mutação -> aptidão (partição de aferição) -> elitismo -> repete; ao fim,
% reavaliação no teste intocado = o envelope reportado.
%
% Uso no Cap. 3 (dentro de um ambiente figure):
%   % ============================================================================
% ESQUEMA PROPOSTO (banca, sugestão — não mergeado): o envelope por algoritmo
% genético (Seção sec:metodo-l0-ag) como laço evolutivo + saída com
% anti-circularidade. A prosa da seção carrega o laço inteiro em um
% parágrafo; a figura o desdobra: população -> torneio -> cruzamento ->
% mutação -> aptidão (partição de aferição) -> elitismo -> repete; ao fim,
% reavaliação no teste intocado = o envelope reportado.
%
% Uso no Cap. 3 (dentro de um ambiente figure):
%   \input{3-metodo/esquemas-propostos/esq-ag-envelope}
%   As citações Holland1975/Goldberg1989 ficam na prosa; a legenda pode
%   remeter ao Apêndice ap:ag (fluxo e operadores formalizados).
% Uso standalone (pré-visualização):
%   \documentclass[tikz,border=6pt]{standalone}
%   \usetikzlibrary{arrows.meta, positioning}
%   \begin{document}\input{esq-ag-envelope.tex}\end{document}
%   (no modo standalone, troque as \ref{...} por texto fixo ou defina-as)
% Requer apenas arrows.meta e positioning (já em packages.tex). Legível em
% P&B: entrada = cinza claro; operação = branco; medição = cinza claro;
% produto = cinza médio; anotação = texto pequeno com tracejado.
% FREEZE respeitado: todos os números são os da própria Seção 3.5.2
% (N_pop=20, 100/200 gerações, k_t=3, p_c=0,8, p_m=0,1, m_s, 10%,
% N_elite=2, 4 cenários, 10 tamanhos, 10 a 30.000), nenhum novo.
% ============================================================================
\begin{tikzpicture}[
    x=1cm, y=1cm,
    caixa/.style={draw, rounded corners=2pt, align=center, text width=34mm,
                  minimum height=11mm, inner sep=3pt, font=\footnotesize},
    entrada/.style={caixa, fill=gray!8},
    medida/.style={caixa, fill=gray!8},
    produto/.style={caixa, fill=gray!14},
    seta/.style={-{Stealth[length=2.4mm]}, thick},
    rot/.style={font=\scriptsize, align=center, inner sep=2pt}
]
    % ------------------ população e proveniência da configuração -----------
    \node[entrada, text width=40mm] (pop) at (0,0)
        {população com $N_{pop}=20$;\\cada indivíduo é um conjunto\\
         de $I$ índices únicos do \textit{pool}};
    \node[rot, anchor=west, text width=54mm, align=left] (prov) at (2.9,-0.1)
        {parâmetros definidos no \textit{notebook}\\
         que executou as corridas (o JSON fixa\\
         apenas $|L_0|$; o padrão $0{,}7$ do cruzamento\\
         é sempre sobrescrito); a população é\\
         o único valor lido do artefato};

    % ------------------ o laço por geração ---------------------------------
    \node[caixa] (torneio) at (0,-2.7)
        {seleção por torneio\\($k_t=3$)};
    \node[caixa] (cruz) at (4.6,-2.7)
        {cruzamento de um ponto\\($p_c=0{,}8$), com reparo\\de unicidade};
    \node[caixa, text width=38mm] (mut) at (9.4,-2.7)
        {mutação (\mbox{$p_m=0{,}1$}):\\
         substitui \mbox{$m_s=\max(1,\lceil 0{,}01\cdot I\rceil)$}\\
         genes};
    \node[medida, text width=38mm] (apt) at (9.4,-5.3)
        {aptidão medida na partição\\de aferição, \textbf{disjunta}\\do
         conjunto de teste};
    \node[caixa] (elit) at (0,-5.3)
        {elitismo de $10\%$\\($N_{elite}=2$): forma\\a nova geração};

    % ------------------ saída: anti-circularidade --------------------------
    \node[produto, text width=52mm] (reaval) at (4.6,-8.2)
        {fim do cenário: os indivíduos finais são \textbf{reavaliados no
         conjunto de teste intocado}; esse valor, tipicamente inferior ao
         de aptidão, é o envelope que os resultados reportam};
    \node[rot, text width=40mm, anchor=north] (circ) at (9.8,-7.4)
        {otimizar e reportar na mesma partição superajustaria o envelope ao
         conjunto de avaliação (circularidade)};

    % ------------------ setas (depois de todos os nós) ----------------------
    \draw[seta] (pop) -- (torneio);
    \draw[seta, dashed] (prov.west) -- (pop.east);
    \draw[seta] (torneio) -- (cruz);
    \draw[seta] (cruz) -- (mut);
    \draw[seta] (mut) -- (apt);
    \draw[seta] (apt) -- (elit);
    \draw[seta] (elit) -- (torneio);
    \draw[seta] (elit.south) |- node[rot, anchor=west, pos=0.28,
        text width=30mm, align=left]
        {após 100 gerações\\(200 só no caso \mbox{$|L_0|=10$})}
        (reaval.west);
    \draw[seta, dashed] (circ.north) -- (apt.south);

    % ------------------ cenários e remissão ---------------------------------
    \node[rot, anchor=west, text width=118mm, align=left] at (-1.9,-10.3)
        {4 cenários por tamanho (maximização e minimização de Acurácia e de
         Macro F1) $\times$ 10 tamanhos de $L_0$, de 10 a 30.000; o fluxo
         completo e os operadores são formalizados no
         Apêndice~\ref{ap:ag}.};
\end{tikzpicture}

%   As citações Holland1975/Goldberg1989 ficam na prosa; a legenda pode
%   remeter ao Apêndice ap:ag (fluxo e operadores formalizados).
% Uso standalone (pré-visualização):
%   \documentclass[tikz,border=6pt]{standalone}
%   \usetikzlibrary{arrows.meta, positioning}
%   \begin{document}% ============================================================================
% ESQUEMA PROPOSTO (banca, sugestão — não mergeado): o envelope por algoritmo
% genético (Seção sec:metodo-l0-ag) como laço evolutivo + saída com
% anti-circularidade. A prosa da seção carrega o laço inteiro em um
% parágrafo; a figura o desdobra: população -> torneio -> cruzamento ->
% mutação -> aptidão (partição de aferição) -> elitismo -> repete; ao fim,
% reavaliação no teste intocado = o envelope reportado.
%
% Uso no Cap. 3 (dentro de um ambiente figure):
%   \input{3-metodo/esquemas-propostos/esq-ag-envelope}
%   As citações Holland1975/Goldberg1989 ficam na prosa; a legenda pode
%   remeter ao Apêndice ap:ag (fluxo e operadores formalizados).
% Uso standalone (pré-visualização):
%   \documentclass[tikz,border=6pt]{standalone}
%   \usetikzlibrary{arrows.meta, positioning}
%   \begin{document}\input{esq-ag-envelope.tex}\end{document}
%   (no modo standalone, troque as \ref{...} por texto fixo ou defina-as)
% Requer apenas arrows.meta e positioning (já em packages.tex). Legível em
% P&B: entrada = cinza claro; operação = branco; medição = cinza claro;
% produto = cinza médio; anotação = texto pequeno com tracejado.
% FREEZE respeitado: todos os números são os da própria Seção 3.5.2
% (N_pop=20, 100/200 gerações, k_t=3, p_c=0,8, p_m=0,1, m_s, 10%,
% N_elite=2, 4 cenários, 10 tamanhos, 10 a 30.000), nenhum novo.
% ============================================================================
\begin{tikzpicture}[
    x=1cm, y=1cm,
    caixa/.style={draw, rounded corners=2pt, align=center, text width=34mm,
                  minimum height=11mm, inner sep=3pt, font=\footnotesize},
    entrada/.style={caixa, fill=gray!8},
    medida/.style={caixa, fill=gray!8},
    produto/.style={caixa, fill=gray!14},
    seta/.style={-{Stealth[length=2.4mm]}, thick},
    rot/.style={font=\scriptsize, align=center, inner sep=2pt}
]
    % ------------------ população e proveniência da configuração -----------
    \node[entrada, text width=40mm] (pop) at (0,0)
        {população com $N_{pop}=20$;\\cada indivíduo é um conjunto\\
         de $I$ índices únicos do \textit{pool}};
    \node[rot, anchor=west, text width=54mm, align=left] (prov) at (2.9,-0.1)
        {parâmetros definidos no \textit{notebook}\\
         que executou as corridas (o JSON fixa\\
         apenas $|L_0|$; o padrão $0{,}7$ do cruzamento\\
         é sempre sobrescrito); a população é\\
         o único valor lido do artefato};

    % ------------------ o laço por geração ---------------------------------
    \node[caixa] (torneio) at (0,-2.7)
        {seleção por torneio\\($k_t=3$)};
    \node[caixa] (cruz) at (4.6,-2.7)
        {cruzamento de um ponto\\($p_c=0{,}8$), com reparo\\de unicidade};
    \node[caixa, text width=38mm] (mut) at (9.4,-2.7)
        {mutação (\mbox{$p_m=0{,}1$}):\\
         substitui \mbox{$m_s=\max(1,\lceil 0{,}01\cdot I\rceil)$}\\
         genes};
    \node[medida, text width=38mm] (apt) at (9.4,-5.3)
        {aptidão medida na partição\\de aferição, \textbf{disjunta}\\do
         conjunto de teste};
    \node[caixa] (elit) at (0,-5.3)
        {elitismo de $10\%$\\($N_{elite}=2$): forma\\a nova geração};

    % ------------------ saída: anti-circularidade --------------------------
    \node[produto, text width=52mm] (reaval) at (4.6,-8.2)
        {fim do cenário: os indivíduos finais são \textbf{reavaliados no
         conjunto de teste intocado}; esse valor, tipicamente inferior ao
         de aptidão, é o envelope que os resultados reportam};
    \node[rot, text width=40mm, anchor=north] (circ) at (9.8,-7.4)
        {otimizar e reportar na mesma partição superajustaria o envelope ao
         conjunto de avaliação (circularidade)};

    % ------------------ setas (depois de todos os nós) ----------------------
    \draw[seta] (pop) -- (torneio);
    \draw[seta, dashed] (prov.west) -- (pop.east);
    \draw[seta] (torneio) -- (cruz);
    \draw[seta] (cruz) -- (mut);
    \draw[seta] (mut) -- (apt);
    \draw[seta] (apt) -- (elit);
    \draw[seta] (elit) -- (torneio);
    \draw[seta] (elit.south) |- node[rot, anchor=west, pos=0.28,
        text width=30mm, align=left]
        {após 100 gerações\\(200 só no caso \mbox{$|L_0|=10$})}
        (reaval.west);
    \draw[seta, dashed] (circ.north) -- (apt.south);

    % ------------------ cenários e remissão ---------------------------------
    \node[rot, anchor=west, text width=118mm, align=left] at (-1.9,-10.3)
        {4 cenários por tamanho (maximização e minimização de Acurácia e de
         Macro F1) $\times$ 10 tamanhos de $L_0$, de 10 a 30.000; o fluxo
         completo e os operadores são formalizados no
         Apêndice~\ref{ap:ag}.};
\end{tikzpicture}
\end{document}
%   (no modo standalone, troque as \ref{...} por texto fixo ou defina-as)
% Requer apenas arrows.meta e positioning (já em packages.tex). Legível em
% P&B: entrada = cinza claro; operação = branco; medição = cinza claro;
% produto = cinza médio; anotação = texto pequeno com tracejado.
% FREEZE respeitado: todos os números são os da própria Seção 3.5.2
% (N_pop=20, 100/200 gerações, k_t=3, p_c=0,8, p_m=0,1, m_s, 10%,
% N_elite=2, 4 cenários, 10 tamanhos, 10 a 30.000), nenhum novo.
% ============================================================================
\begin{tikzpicture}[
    x=1cm, y=1cm,
    caixa/.style={draw, rounded corners=2pt, align=center, text width=34mm,
                  minimum height=11mm, inner sep=3pt, font=\footnotesize},
    entrada/.style={caixa, fill=gray!8},
    medida/.style={caixa, fill=gray!8},
    produto/.style={caixa, fill=gray!14},
    seta/.style={-{Stealth[length=2.4mm]}, thick},
    rot/.style={font=\scriptsize, align=center, inner sep=2pt}
]
    % ------------------ população e proveniência da configuração -----------
    \node[entrada, text width=40mm] (pop) at (0,0)
        {população com $N_{pop}=20$;\\cada indivíduo é um conjunto\\
         de $I$ índices únicos do \textit{pool}};
    \node[rot, anchor=west, text width=54mm, align=left] (prov) at (2.9,-0.1)
        {parâmetros definidos no \textit{notebook}\\
         que executou as corridas (o JSON fixa\\
         apenas $|L_0|$; o padrão $0{,}7$ do cruzamento\\
         é sempre sobrescrito); a população é\\
         o único valor lido do artefato};

    % ------------------ o laço por geração ---------------------------------
    \node[caixa] (torneio) at (0,-2.7)
        {seleção por torneio\\($k_t=3$)};
    \node[caixa] (cruz) at (4.6,-2.7)
        {cruzamento de um ponto\\($p_c=0{,}8$), com reparo\\de unicidade};
    \node[caixa, text width=38mm] (mut) at (9.4,-2.7)
        {mutação (\mbox{$p_m=0{,}1$}):\\
         substitui \mbox{$m_s=\max(1,\lceil 0{,}01\cdot I\rceil)$}\\
         genes};
    \node[medida, text width=38mm] (apt) at (9.4,-5.3)
        {aptidão medida na partição\\de aferição, \textbf{disjunta}\\do
         conjunto de teste};
    \node[caixa] (elit) at (0,-5.3)
        {elitismo de $10\%$\\($N_{elite}=2$): forma\\a nova geração};

    % ------------------ saída: anti-circularidade --------------------------
    \node[produto, text width=52mm] (reaval) at (4.6,-8.2)
        {fim do cenário: os indivíduos finais são \textbf{reavaliados no
         conjunto de teste intocado}; esse valor, tipicamente inferior ao
         de aptidão, é o envelope que os resultados reportam};
    \node[rot, text width=40mm, anchor=north] (circ) at (9.8,-7.4)
        {otimizar e reportar na mesma partição superajustaria o envelope ao
         conjunto de avaliação (circularidade)};

    % ------------------ setas (depois de todos os nós) ----------------------
    \draw[seta] (pop) -- (torneio);
    \draw[seta, dashed] (prov.west) -- (pop.east);
    \draw[seta] (torneio) -- (cruz);
    \draw[seta] (cruz) -- (mut);
    \draw[seta] (mut) -- (apt);
    \draw[seta] (apt) -- (elit);
    \draw[seta] (elit) -- (torneio);
    \draw[seta] (elit.south) |- node[rot, anchor=west, pos=0.28,
        text width=30mm, align=left]
        {após 100 gerações\\(200 só no caso \mbox{$|L_0|=10$})}
        (reaval.west);
    \draw[seta, dashed] (circ.north) -- (apt.south);

    % ------------------ cenários e remissão ---------------------------------
    \node[rot, anchor=west, text width=118mm, align=left] at (-1.9,-10.3)
        {4 cenários por tamanho (maximização e minimização de Acurácia e de
         Macro F1) $\times$ 10 tamanhos de $L_0$, de 10 a 30.000; o fluxo
         completo e os operadores são formalizados no
         Apêndice~\ref{ap:ag}.};
\end{tikzpicture}

%   As citações Holland1975/Goldberg1989 ficam na prosa; a legenda pode
%   remeter ao Apêndice ap:ag (fluxo e operadores formalizados).
% Uso standalone (pré-visualização):
%   \documentclass[tikz,border=6pt]{standalone}
%   \usetikzlibrary{arrows.meta, positioning}
%   \begin{document}% ============================================================================
% ESQUEMA PROPOSTO (banca, sugestão — não mergeado): o envelope por algoritmo
% genético (Seção sec:metodo-l0-ag) como laço evolutivo + saída com
% anti-circularidade. A prosa da seção carrega o laço inteiro em um
% parágrafo; a figura o desdobra: população -> torneio -> cruzamento ->
% mutação -> aptidão (partição de aferição) -> elitismo -> repete; ao fim,
% reavaliação no teste intocado = o envelope reportado.
%
% Uso no Cap. 3 (dentro de um ambiente figure):
%   % ============================================================================
% ESQUEMA PROPOSTO (banca, sugestão — não mergeado): o envelope por algoritmo
% genético (Seção sec:metodo-l0-ag) como laço evolutivo + saída com
% anti-circularidade. A prosa da seção carrega o laço inteiro em um
% parágrafo; a figura o desdobra: população -> torneio -> cruzamento ->
% mutação -> aptidão (partição de aferição) -> elitismo -> repete; ao fim,
% reavaliação no teste intocado = o envelope reportado.
%
% Uso no Cap. 3 (dentro de um ambiente figure):
%   \input{3-metodo/esquemas-propostos/esq-ag-envelope}
%   As citações Holland1975/Goldberg1989 ficam na prosa; a legenda pode
%   remeter ao Apêndice ap:ag (fluxo e operadores formalizados).
% Uso standalone (pré-visualização):
%   \documentclass[tikz,border=6pt]{standalone}
%   \usetikzlibrary{arrows.meta, positioning}
%   \begin{document}\input{esq-ag-envelope.tex}\end{document}
%   (no modo standalone, troque as \ref{...} por texto fixo ou defina-as)
% Requer apenas arrows.meta e positioning (já em packages.tex). Legível em
% P&B: entrada = cinza claro; operação = branco; medição = cinza claro;
% produto = cinza médio; anotação = texto pequeno com tracejado.
% FREEZE respeitado: todos os números são os da própria Seção 3.5.2
% (N_pop=20, 100/200 gerações, k_t=3, p_c=0,8, p_m=0,1, m_s, 10%,
% N_elite=2, 4 cenários, 10 tamanhos, 10 a 30.000), nenhum novo.
% ============================================================================
\begin{tikzpicture}[
    x=1cm, y=1cm,
    caixa/.style={draw, rounded corners=2pt, align=center, text width=34mm,
                  minimum height=11mm, inner sep=3pt, font=\footnotesize},
    entrada/.style={caixa, fill=gray!8},
    medida/.style={caixa, fill=gray!8},
    produto/.style={caixa, fill=gray!14},
    seta/.style={-{Stealth[length=2.4mm]}, thick},
    rot/.style={font=\scriptsize, align=center, inner sep=2pt}
]
    % ------------------ população e proveniência da configuração -----------
    \node[entrada, text width=40mm] (pop) at (0,0)
        {população com $N_{pop}=20$;\\cada indivíduo é um conjunto\\
         de $I$ índices únicos do \textit{pool}};
    \node[rot, anchor=west, text width=54mm, align=left] (prov) at (2.9,-0.1)
        {parâmetros definidos no \textit{notebook}\\
         que executou as corridas (o JSON fixa\\
         apenas $|L_0|$; o padrão $0{,}7$ do cruzamento\\
         é sempre sobrescrito); a população é\\
         o único valor lido do artefato};

    % ------------------ o laço por geração ---------------------------------
    \node[caixa] (torneio) at (0,-2.7)
        {seleção por torneio\\($k_t=3$)};
    \node[caixa] (cruz) at (4.6,-2.7)
        {cruzamento de um ponto\\($p_c=0{,}8$), com reparo\\de unicidade};
    \node[caixa, text width=38mm] (mut) at (9.4,-2.7)
        {mutação (\mbox{$p_m=0{,}1$}):\\
         substitui \mbox{$m_s=\max(1,\lceil 0{,}01\cdot I\rceil)$}\\
         genes};
    \node[medida, text width=38mm] (apt) at (9.4,-5.3)
        {aptidão medida na partição\\de aferição, \textbf{disjunta}\\do
         conjunto de teste};
    \node[caixa] (elit) at (0,-5.3)
        {elitismo de $10\%$\\($N_{elite}=2$): forma\\a nova geração};

    % ------------------ saída: anti-circularidade --------------------------
    \node[produto, text width=52mm] (reaval) at (4.6,-8.2)
        {fim do cenário: os indivíduos finais são \textbf{reavaliados no
         conjunto de teste intocado}; esse valor, tipicamente inferior ao
         de aptidão, é o envelope que os resultados reportam};
    \node[rot, text width=40mm, anchor=north] (circ) at (9.8,-7.4)
        {otimizar e reportar na mesma partição superajustaria o envelope ao
         conjunto de avaliação (circularidade)};

    % ------------------ setas (depois de todos os nós) ----------------------
    \draw[seta] (pop) -- (torneio);
    \draw[seta, dashed] (prov.west) -- (pop.east);
    \draw[seta] (torneio) -- (cruz);
    \draw[seta] (cruz) -- (mut);
    \draw[seta] (mut) -- (apt);
    \draw[seta] (apt) -- (elit);
    \draw[seta] (elit) -- (torneio);
    \draw[seta] (elit.south) |- node[rot, anchor=west, pos=0.28,
        text width=30mm, align=left]
        {após 100 gerações\\(200 só no caso \mbox{$|L_0|=10$})}
        (reaval.west);
    \draw[seta, dashed] (circ.north) -- (apt.south);

    % ------------------ cenários e remissão ---------------------------------
    \node[rot, anchor=west, text width=118mm, align=left] at (-1.9,-10.3)
        {4 cenários por tamanho (maximização e minimização de Acurácia e de
         Macro F1) $\times$ 10 tamanhos de $L_0$, de 10 a 30.000; o fluxo
         completo e os operadores são formalizados no
         Apêndice~\ref{ap:ag}.};
\end{tikzpicture}

%   As citações Holland1975/Goldberg1989 ficam na prosa; a legenda pode
%   remeter ao Apêndice ap:ag (fluxo e operadores formalizados).
% Uso standalone (pré-visualização):
%   \documentclass[tikz,border=6pt]{standalone}
%   \usetikzlibrary{arrows.meta, positioning}
%   \begin{document}% ============================================================================
% ESQUEMA PROPOSTO (banca, sugestão — não mergeado): o envelope por algoritmo
% genético (Seção sec:metodo-l0-ag) como laço evolutivo + saída com
% anti-circularidade. A prosa da seção carrega o laço inteiro em um
% parágrafo; a figura o desdobra: população -> torneio -> cruzamento ->
% mutação -> aptidão (partição de aferição) -> elitismo -> repete; ao fim,
% reavaliação no teste intocado = o envelope reportado.
%
% Uso no Cap. 3 (dentro de um ambiente figure):
%   \input{3-metodo/esquemas-propostos/esq-ag-envelope}
%   As citações Holland1975/Goldberg1989 ficam na prosa; a legenda pode
%   remeter ao Apêndice ap:ag (fluxo e operadores formalizados).
% Uso standalone (pré-visualização):
%   \documentclass[tikz,border=6pt]{standalone}
%   \usetikzlibrary{arrows.meta, positioning}
%   \begin{document}\input{esq-ag-envelope.tex}\end{document}
%   (no modo standalone, troque as \ref{...} por texto fixo ou defina-as)
% Requer apenas arrows.meta e positioning (já em packages.tex). Legível em
% P&B: entrada = cinza claro; operação = branco; medição = cinza claro;
% produto = cinza médio; anotação = texto pequeno com tracejado.
% FREEZE respeitado: todos os números são os da própria Seção 3.5.2
% (N_pop=20, 100/200 gerações, k_t=3, p_c=0,8, p_m=0,1, m_s, 10%,
% N_elite=2, 4 cenários, 10 tamanhos, 10 a 30.000), nenhum novo.
% ============================================================================
\begin{tikzpicture}[
    x=1cm, y=1cm,
    caixa/.style={draw, rounded corners=2pt, align=center, text width=34mm,
                  minimum height=11mm, inner sep=3pt, font=\footnotesize},
    entrada/.style={caixa, fill=gray!8},
    medida/.style={caixa, fill=gray!8},
    produto/.style={caixa, fill=gray!14},
    seta/.style={-{Stealth[length=2.4mm]}, thick},
    rot/.style={font=\scriptsize, align=center, inner sep=2pt}
]
    % ------------------ população e proveniência da configuração -----------
    \node[entrada, text width=40mm] (pop) at (0,0)
        {população com $N_{pop}=20$;\\cada indivíduo é um conjunto\\
         de $I$ índices únicos do \textit{pool}};
    \node[rot, anchor=west, text width=54mm, align=left] (prov) at (2.9,-0.1)
        {parâmetros definidos no \textit{notebook}\\
         que executou as corridas (o JSON fixa\\
         apenas $|L_0|$; o padrão $0{,}7$ do cruzamento\\
         é sempre sobrescrito); a população é\\
         o único valor lido do artefato};

    % ------------------ o laço por geração ---------------------------------
    \node[caixa] (torneio) at (0,-2.7)
        {seleção por torneio\\($k_t=3$)};
    \node[caixa] (cruz) at (4.6,-2.7)
        {cruzamento de um ponto\\($p_c=0{,}8$), com reparo\\de unicidade};
    \node[caixa, text width=38mm] (mut) at (9.4,-2.7)
        {mutação (\mbox{$p_m=0{,}1$}):\\
         substitui \mbox{$m_s=\max(1,\lceil 0{,}01\cdot I\rceil)$}\\
         genes};
    \node[medida, text width=38mm] (apt) at (9.4,-5.3)
        {aptidão medida na partição\\de aferição, \textbf{disjunta}\\do
         conjunto de teste};
    \node[caixa] (elit) at (0,-5.3)
        {elitismo de $10\%$\\($N_{elite}=2$): forma\\a nova geração};

    % ------------------ saída: anti-circularidade --------------------------
    \node[produto, text width=52mm] (reaval) at (4.6,-8.2)
        {fim do cenário: os indivíduos finais são \textbf{reavaliados no
         conjunto de teste intocado}; esse valor, tipicamente inferior ao
         de aptidão, é o envelope que os resultados reportam};
    \node[rot, text width=40mm, anchor=north] (circ) at (9.8,-7.4)
        {otimizar e reportar na mesma partição superajustaria o envelope ao
         conjunto de avaliação (circularidade)};

    % ------------------ setas (depois de todos os nós) ----------------------
    \draw[seta] (pop) -- (torneio);
    \draw[seta, dashed] (prov.west) -- (pop.east);
    \draw[seta] (torneio) -- (cruz);
    \draw[seta] (cruz) -- (mut);
    \draw[seta] (mut) -- (apt);
    \draw[seta] (apt) -- (elit);
    \draw[seta] (elit) -- (torneio);
    \draw[seta] (elit.south) |- node[rot, anchor=west, pos=0.28,
        text width=30mm, align=left]
        {após 100 gerações\\(200 só no caso \mbox{$|L_0|=10$})}
        (reaval.west);
    \draw[seta, dashed] (circ.north) -- (apt.south);

    % ------------------ cenários e remissão ---------------------------------
    \node[rot, anchor=west, text width=118mm, align=left] at (-1.9,-10.3)
        {4 cenários por tamanho (maximização e minimização de Acurácia e de
         Macro F1) $\times$ 10 tamanhos de $L_0$, de 10 a 30.000; o fluxo
         completo e os operadores são formalizados no
         Apêndice~\ref{ap:ag}.};
\end{tikzpicture}
\end{document}
%   (no modo standalone, troque as \ref{...} por texto fixo ou defina-as)
% Requer apenas arrows.meta e positioning (já em packages.tex). Legível em
% P&B: entrada = cinza claro; operação = branco; medição = cinza claro;
% produto = cinza médio; anotação = texto pequeno com tracejado.
% FREEZE respeitado: todos os números são os da própria Seção 3.5.2
% (N_pop=20, 100/200 gerações, k_t=3, p_c=0,8, p_m=0,1, m_s, 10%,
% N_elite=2, 4 cenários, 10 tamanhos, 10 a 30.000), nenhum novo.
% ============================================================================
\begin{tikzpicture}[
    x=1cm, y=1cm,
    caixa/.style={draw, rounded corners=2pt, align=center, text width=34mm,
                  minimum height=11mm, inner sep=3pt, font=\footnotesize},
    entrada/.style={caixa, fill=gray!8},
    medida/.style={caixa, fill=gray!8},
    produto/.style={caixa, fill=gray!14},
    seta/.style={-{Stealth[length=2.4mm]}, thick},
    rot/.style={font=\scriptsize, align=center, inner sep=2pt}
]
    % ------------------ população e proveniência da configuração -----------
    \node[entrada, text width=40mm] (pop) at (0,0)
        {população com $N_{pop}=20$;\\cada indivíduo é um conjunto\\
         de $I$ índices únicos do \textit{pool}};
    \node[rot, anchor=west, text width=54mm, align=left] (prov) at (2.9,-0.1)
        {parâmetros definidos no \textit{notebook}\\
         que executou as corridas (o JSON fixa\\
         apenas $|L_0|$; o padrão $0{,}7$ do cruzamento\\
         é sempre sobrescrito); a população é\\
         o único valor lido do artefato};

    % ------------------ o laço por geração ---------------------------------
    \node[caixa] (torneio) at (0,-2.7)
        {seleção por torneio\\($k_t=3$)};
    \node[caixa] (cruz) at (4.6,-2.7)
        {cruzamento de um ponto\\($p_c=0{,}8$), com reparo\\de unicidade};
    \node[caixa, text width=38mm] (mut) at (9.4,-2.7)
        {mutação (\mbox{$p_m=0{,}1$}):\\
         substitui \mbox{$m_s=\max(1,\lceil 0{,}01\cdot I\rceil)$}\\
         genes};
    \node[medida, text width=38mm] (apt) at (9.4,-5.3)
        {aptidão medida na partição\\de aferição, \textbf{disjunta}\\do
         conjunto de teste};
    \node[caixa] (elit) at (0,-5.3)
        {elitismo de $10\%$\\($N_{elite}=2$): forma\\a nova geração};

    % ------------------ saída: anti-circularidade --------------------------
    \node[produto, text width=52mm] (reaval) at (4.6,-8.2)
        {fim do cenário: os indivíduos finais são \textbf{reavaliados no
         conjunto de teste intocado}; esse valor, tipicamente inferior ao
         de aptidão, é o envelope que os resultados reportam};
    \node[rot, text width=40mm, anchor=north] (circ) at (9.8,-7.4)
        {otimizar e reportar na mesma partição superajustaria o envelope ao
         conjunto de avaliação (circularidade)};

    % ------------------ setas (depois de todos os nós) ----------------------
    \draw[seta] (pop) -- (torneio);
    \draw[seta, dashed] (prov.west) -- (pop.east);
    \draw[seta] (torneio) -- (cruz);
    \draw[seta] (cruz) -- (mut);
    \draw[seta] (mut) -- (apt);
    \draw[seta] (apt) -- (elit);
    \draw[seta] (elit) -- (torneio);
    \draw[seta] (elit.south) |- node[rot, anchor=west, pos=0.28,
        text width=30mm, align=left]
        {após 100 gerações\\(200 só no caso \mbox{$|L_0|=10$})}
        (reaval.west);
    \draw[seta, dashed] (circ.north) -- (apt.south);

    % ------------------ cenários e remissão ---------------------------------
    \node[rot, anchor=west, text width=118mm, align=left] at (-1.9,-10.3)
        {4 cenários por tamanho (maximização e minimização de Acurácia e de
         Macro F1) $\times$ 10 tamanhos de $L_0$, de 10 a 30.000; o fluxo
         completo e os operadores são formalizados no
         Apêndice~\ref{ap:ag}.};
\end{tikzpicture}
\end{document}
%   (no modo standalone, troque as \ref{...} por texto fixo ou defina-as)
% Requer apenas arrows.meta e positioning (já em packages.tex). Legível em
% P&B: entrada = cinza claro; operação = branco; medição = cinza claro;
% produto = cinza médio; anotação = texto pequeno com tracejado.
% FREEZE respeitado: todos os números são os da própria Seção 3.5.2
% (N_pop=20, 100/200 gerações, k_t=3, p_c=0,8, p_m=0,1, m_s, 10%,
% N_elite=2, 4 cenários, 10 tamanhos, 10 a 30.000), nenhum novo.
% ============================================================================
\begin{tikzpicture}[
    x=1cm, y=1cm,
    caixa/.style={draw, rounded corners=2pt, align=center, text width=34mm,
                  minimum height=11mm, inner sep=3pt, font=\footnotesize},
    entrada/.style={caixa, fill=gray!8},
    medida/.style={caixa, fill=gray!8},
    produto/.style={caixa, fill=gray!14},
    seta/.style={-{Stealth[length=2.4mm]}, thick},
    rot/.style={font=\scriptsize, align=center, inner sep=2pt}
]
    % ------------------ população e proveniência da configuração -----------
    \node[entrada, text width=40mm] (pop) at (0,0)
        {população com $N_{pop}=20$;\\cada indivíduo é um conjunto\\
         de $I$ índices únicos do \textit{pool}};
    \node[rot, anchor=west, text width=54mm, align=left] (prov) at (2.9,-0.1)
        {parâmetros definidos no \textit{notebook}\\
         que executou as corridas (o JSON fixa\\
         apenas $|L_0|$; o padrão $0{,}7$ do cruzamento\\
         é sempre sobrescrito); a população é\\
         o único valor lido do artefato};

    % ------------------ o laço por geração ---------------------------------
    \node[caixa] (torneio) at (0,-2.7)
        {seleção por torneio\\($k_t=3$)};
    \node[caixa] (cruz) at (4.6,-2.7)
        {cruzamento de um ponto\\($p_c=0{,}8$), com reparo\\de unicidade};
    \node[caixa, text width=38mm] (mut) at (9.4,-2.7)
        {mutação (\mbox{$p_m=0{,}1$}):\\
         substitui \mbox{$m_s=\max(1,\lceil 0{,}01\cdot I\rceil)$}\\
         genes};
    \node[medida, text width=38mm] (apt) at (9.4,-5.3)
        {aptidão medida na partição\\de aferição, \textbf{disjunta}\\do
         conjunto de teste};
    \node[caixa] (elit) at (0,-5.3)
        {elitismo de $10\%$\\($N_{elite}=2$): forma\\a nova geração};

    % ------------------ saída: anti-circularidade --------------------------
    \node[produto, text width=52mm] (reaval) at (4.6,-8.2)
        {fim do cenário: os indivíduos finais são \textbf{reavaliados no
         conjunto de teste intocado}; esse valor, tipicamente inferior ao
         de aptidão, é o envelope que os resultados reportam};
    \node[rot, text width=40mm, anchor=north] (circ) at (9.8,-7.4)
        {otimizar e reportar na mesma partição superajustaria o envelope ao
         conjunto de avaliação (circularidade)};

    % ------------------ setas (depois de todos os nós) ----------------------
    \draw[seta] (pop) -- (torneio);
    \draw[seta, dashed] (prov.west) -- (pop.east);
    \draw[seta] (torneio) -- (cruz);
    \draw[seta] (cruz) -- (mut);
    \draw[seta] (mut) -- (apt);
    \draw[seta] (apt) -- (elit);
    \draw[seta] (elit) -- (torneio);
    \draw[seta] (elit.south) |- node[rot, anchor=west, pos=0.28,
        text width=30mm, align=left]
        {após 100 gerações\\(200 só no caso \mbox{$|L_0|=10$})}
        (reaval.west);
    \draw[seta, dashed] (circ.north) -- (apt.south);

    % ------------------ cenários e remissão ---------------------------------
    \node[rot, anchor=west, text width=118mm, align=left] at (-1.9,-10.3)
        {4 cenários por tamanho (maximização e minimização de Acurácia e de
         Macro F1) $\times$ 10 tamanhos de $L_0$, de 10 a 30.000; o fluxo
         completo e os operadores são formalizados no
         Apêndice~\ref{ap:ag}.};
\end{tikzpicture}

    \caption{O laço do algoritmo genético e a saída anticircular, em resumo
    visual: a aptidão que orienta a busca é medida na partição de aferição,
    disjunta do teste, e o envelope reportado é o do melhor indivíduo
    reavaliado no conjunto de teste intocado.}
    \label{fig:ap-ag-laco}
\end{figure}

A implementação de referência da reexecução e seus resultados acompanham o
repositório de código da tese.
